\subsection{U1$\to$U2 falsifiable interface chains (benchmark consequences without theorem-level upgrade)}
\label{app:u1_to_u2_falsifiable_interface_chains}

This appendix decomposes the informal statement ``a simple-group unification (U1) would explain coupling relations (U2)'' into a small set of auditable, falsifiable interface chains.
No U1 theorem is claimed; the goal is to expose the minimal failure points that prevent upgrading U1 in the current tick+CAP closure discipline.

\subsubsection{Chain C1: embedding-fixes-normalization $\Rightarrow$ constrained U2 comparison}
\paragraph{Statement (interface).}
\InterfaceTag If a simple group $G$ embeds the SM factors and fixes the normalization of the abelian generator, then the conversion between $\alpha_Y$ and the benchmark $\alpha_1$ is not arbitrary; it is fixed by the embedding (Appendix~\ref{app:u3_normalization_embedding_registry}).
Therefore a U2 audit must be evaluated under the embedding-fixed U3 convention, not under a free rescaling.

\paragraph{Minimal failure points.}
\AuditTag The upgrade C1$\Rightarrow$U1 fails (within this manuscript) if any of the following holds:
\begin{itemize}
\item (C1-F1) the embedding-fixing of U3 is not supplied as a closed input (it is currently only a bounded registry in Appendix~\ref{app:u3_normalization_embedding_registry});
\item (C1-F2) the U2 audit output depends strongly on U3 convention choice, so that U1 cannot be inferred without an independent convention-fixing channel;
\item (C1-F3) threshold/scheme uncertainties (Appendix~\ref{app:coupling_unification_audit_in_r}) dominate the intersection mismatch so that any ``unification'' reading is non-robust.
\end{itemize}

\subsubsection{Chain C2: U1 predicts correlated threshold structure (pattern constraint)}
\paragraph{Statement (interface).}
\InterfaceTag In a genuine U1 completion, heavy degrees of freedom at or near a unification scale introduce correlated threshold corrections across $(g_1,g_2,g_3)$ rather than independent per-factor shifts.
As an interface consequence, one may register a bounded family of correlated threshold-pattern hypotheses and compare them to the observed U2 mismatch pattern.

\paragraph{Minimal failure points.}
\AuditTag The upgrade C2$\Rightarrow$U1 fails (within this manuscript) if:
\begin{itemize}
\item (C2-F1) no bounded correlated threshold-pattern family is supplied that is not tuned to match U2;
\item (C2-F2) the observed mismatch pattern can be equally well explained by per-factor scheme shifts, making the correlation non-identifying;
\item (C2-F3) the required correction pattern would force denominators or parameter bounds far beyond the paper's CAP-minimal audit families.
\end{itemize}

\subsubsection{Chain C3: U1 implies additional high-energy channels beyond U2 (independent falsifiers)}
\paragraph{Statement (interface).}
\InterfaceTag U2 alone is insufficient to upgrade U1: a simple-group completion typically implies additional observable channels (e.g., baryon-number violation, monopoles, or specific representation constraints).
Within this paper's discipline, any U1 upgrade requires at least one independent falsifier chain that does not reuse the U2 matching target.

\paragraph{Minimal failure points.}
\AuditTag The upgrade C3$\Rightarrow$U1 fails (within this manuscript) if:
\begin{itemize}
\item (C3-F1) the additional channels are not connected to a vendored dataset and a deterministic audit pipeline;
\item (C3-F2) the predicted channel rates/structures depend on model-building inputs not present in the tick+CAP closure chain;
\item (C3-F3) the bridge requires scattering/renormalization claims beyond the scoped interfaces (Appendix~\ref{app:scattering_haag_ruelle_lsz_interface}; Appendix~\ref{app:renormalization_dictionary_and_boundaries}).
\end{itemize}
